\chapter{}

\iffalse
Gibbs自由能$G(T,P,N)=\mu N$，因此化学势是$T,P$的函数。一种测量化学势的方法是利用吸附现象。暴露在气体或液体中的表面，在某些位置上会吸附气体或液体中的分子。吸附和脱附过程是重要的表面现象。考虑最简单的情形，吸附能$\varepsilon<0$，统计力学给出该吸附位占据的几率是
\[
	p=\frac1{1+\e^{(\varepsilon-\mu)/\kB T}}.
\]
\begin{notecaptionblock}
	\tikzfigure{adsorption-occupation}
	\captionof{figure}{吸附位占据几率}
\end{notecaptionblock}
因此
\[
	\mu=\varepsilon-\kB T\ln\bigp{p^{-1}-1}.
\]

\sectionstar{非平衡热力学}
\subsection*{输运现象}
热导中，热流密度
\[\bm j_q=-\kappa\,\grad T.\]
其中$\kappa$为热导率。
	电导中，电流密度
\[\bm j_e=-\sigma\,\grad\phi.\]
即Ohm定律，其中$\sigma$为电导率。
	扩散中，
\[\bm j_n=D\,\grad n.\]

\subsection*{熵的产生}
熵流
\[\dv st=\pv st+\div \bm J_s\]
这样热流密度
\[\bm J_q=T\bm J_s=\bm J_u-\mu\bm J_n\]
非平衡热力学基本方程
\[\dv st=\sum\bm F_i\cdot\bm J_i.\]
其中$\bm F_i:=\grad A_i$为驱动力。

\newcommand*{\entransy}{\scalebox{.4}[1]{火}\hspace{-4.3pt}\scalebox{.4}[1]{火}\hspace{-1pt}\scalebox{.7}[1]{积}}

\subsection*{\entransy}
\entransy\footnote{笑死，这个字根本打不出来。}(entransy)是清华大学过增元院士等通过热电类比提出的一个新的概念。顾名思义，\entransy 被定义为热能和温度的积（的一半）
\[G:=\frac12QT.\]
在过院士等人的著作《传热\entransy 理论及其应用》中，\entransy 有热势能的含义，并类比熵的英文entropy，将\entransy 的英文称作entransy。在书中，作者基于能量守恒方程推导了平衡方程，并提出耗散和耗散热阻等概念。通过平衡方程建立了系统内部传热引起的耗散与边界换热温差、热流之间的关系，为传热过程和热系统的优化提供了新的方法。%\underline{并差点得了科技进步奖。}
感兴趣的同学春季学期可以选修\textit{新概念热学(New Thermology)}\footnote{某老师更愿意简称其为NT。}。

\fi